\section{绪\hspace{1em}论}

本章首先介绍智能体工作流自动生成领域的研究背景与发展趋势，分析当前方法在专业领域知识利用和智能体匹配方面存在的核心挑战，随后综述国内外相关研究现状，最后阐明本研究的主要内容与论文结构。

\subsection{课题背景、目的与意义}

\subsubsection{研究背景与趋势}

近年来，以GPT-4、Claude\footnote{Anthropic. Claude 3 Model Card. \url{https://assets.anthropic.com/m/61e7d27f8c8f5919/original/Claude-3-Model-Card.pdf}}、DeepSeek为代表的大型语言模型（Large Language Model，LLM）在自然语言理解、逻辑推理与代码生成等任务上取得了显著突破\cite{openai2023gpt4}，国内研究也开始系统总结大语言模型评测方法与应用风险\cite{luo2024llmeval}。在此基础上，研究者们将LLM的能力进一步延伸至智能体（Agent）系统，通过赋予模型工具调用、记忆管理与多步规划能力，构建能够自主完成复杂任务的智能体框架。

多智能体协作（Multi-Agent Collaboration）是当前智能体研究的重要方向。与单一智能体相比，多智能体系统通过角色分工与并行协作，能够处理更加复杂的长链路任务。LangGraph\footnote{LangChain AI. LangGraph: Building Stateful Multi-Actor Applications with LLMs. \url{https://github.com/langchain-ai/langgraph}}、AutoGen\cite{wu2023autogen}等框架相继涌现，为多智能体工作流的编程式定义提供了基础支撑。然而，上述框架均依赖人工预先设计工作流拓扑，难以应对开放领域的动态任务需求。

与此同时，检索增强生成（Retrieval-Augmented Generation，RAG）技术通过将外部知识注入LLM的推理过程，有效缓解了模型幻觉问题，并显著提升了领域适配能力\cite{lewis2020retrieval}。进一步地，以Microsoft GraphRAG\cite{edge2024graphrag}为代表的图结构检索方法，将知识组织为实体-关系图，利用图遍历实现多跳推理，在复杂问答和摘要任务上展现出超越传统向量检索的性能。

电力系统作为关系国计民生的重要基础设施，其运维过程涉及大量专业知识与复杂操作规程\footnote{NERC. Lessons Learned: Protective Relay System Misoperations. \url{https://www.nerc.com/pa/rrm/ea/Lessons\%20Learned\%20Document\%20Library}}。电力设备故障排查、保护定值校验、新能源消纳分析等任务，需要多个专业角色协同完成，天然契合多智能体工作流的应用场景。已有研究利用深度学习与知识图谱进行变电站设备故障诊断\cite{xiao2022substationkg}，也有研究将知识图谱与溯因推理结合用于智能电网故障诊断\cite{yao2024smartgridkg}。此外，电力系统存在以CIM（Common Information Model）为代表的标准化知识本体，适合以知识图谱的形式进行结构化表达，为GraphRAG方法的应用提供了优质的数据基础。

\subsubsection{面临的问题与挑战}

尽管LLM在通用任务上表现出色，将其直接用于电力等专业领域的工作流自动生成时，仍面临以下核心挑战：

\textbf{（1）领域知识获取不足。}LLM的参数化知识主要来源于互联网通用语料，对专业领域内的最新标准规程、设备配置文档和故障案例覆盖有限。在缺乏外部知识注入的情况下，生成的工作流往往缺乏专业性和可操作性。

\textbf{（2）多模态知识融合困难。}电力系统文档普遍包含文本、表格、单线图和流程图等多种信息形态。现有的RAG方法大多仅处理纯文本，难以有效利用图像和图形所蕴含的结构化信息，导致知识抽取不完整。

\textbf{（3）智能体匹配缺乏系统化机制。}现有多智能体框架通常依赖LLM自身的理解能力进行隐式的角色分配，缺乏对智能体能力（Capability）、工具（Tool）和接口规范（Schema）的显式建模与量化匹配，造成角色分配的随机性与低准确率。

\textbf{（4）工作流质量难以保证。}自动生成的工作流可能存在循环依赖、数据流断裂等结构性错误，以及逻辑顺序不合理等语义性问题，缺乏系统性的验证机制。

\subsubsection{课题目的与意义}

针对上述挑战，本研究旨在构建一套完整的多模态GraphRAG驱动的智能体工作流自主生成方法，实现对专业领域任务的准确分解、知识增强和角色分配。具体目标如下：

一是构建面向电力领域的多模态知识图谱，支持文本、图像与流程图的统一表示和向量检索；二是设计基于子图扩展的GraphRAG检索方法，并定量分析向量种子检索与多跳邻域扩展在下游任务中的独立贡献；三是提出量化的子任务-智能体匹配算法，从能力、工具、接口三个维度对智能体进行评分与优选；四是实现工作流的DAG编排与迭代验证-修复机制，保证生成工作流的可执行性和逻辑合理性。

本研究的成果可为电力调度、故障应急处置、知识服务等场景提供可靠的智能体工作流自动化支撑，具有重要的工程应用价值；同时，所提方法具有良好的领域通用性，可推广至医疗、法律、金融等其他专业领域。

\subsection{国内外研究现状}

\subsubsection{检索增强生成方法}

检索增强生成（RAG）由Lewis等人\cite{lewis2020retrieval}于2020年正式提出，通过在推理阶段动态检索外部文档并将其作为上下文输入LLM，显著缓解了模型的事实幻觉问题。早期RAG方法基于稠密向量检索（Dense Passage Retrieval，DPR），以FAISS等近似最近邻索引为核心，在开放域问答任务上取得了良好效果。

随后，研究者从检索粒度、融合策略和迭代机制等多个维度对RAG进行了拓展。Self-RAG\cite{asai2023selfrag}引入反思机制，使模型能够动态判断是否需要检索；CRAG\cite{yan2024crag}提出上下文自适应压缩策略，提升了长文档场景下的检索效率；Adaptive RAG\cite{jeong2024adaptive}根据问题复杂度自适应选择检索策略。上述方法均以向量相似度为核心度量，在多跳推理和结构化知识利用方面存在固有局限。

\subsubsection{基于图结构的知识检索}

知识图谱（Knowledge Graph，KG）为结构化知识的存储和推理提供了天然载体，国内学者已对知识图谱的知识抽取、表示、融合和推理技术进行系统梳理\cite{xu2016kgreview}，并对知识表示学习的发展进行了总结\cite{liu2016krl}。KGRAG\cite{kgrag2023}将传统RAG与知识图谱结合，通过实体链接将检索问题转化为子图查询，在复杂关系推理场景下显著优于向量检索。GraphRAG\cite{edge2024graphrag}由微软研究院提出，在构建图的基础上引入社区摘要（Community Summary），实现了从局部到全局的层次化知识表示，在长文档摘要任务上展现出卓越性能。

RAPTOR\cite{sarthi2024raptor}通过递归文本摘要树构建多层次上下文索引，G-RAG\cite{grag2024}引入图神经网络对检索节点进行语义重排。近期，HippoRAG\cite{gutierrez2024hipporag}借鉴人类海马体记忆机制，构建个性化知识图谱以支持长期记忆检索。国内关于知识图谱与大语言模型融合的综述也指出，结构化知识有助于缓解幻觉并提升可解释性\cite{cao2025kgllmreview}。然而，上述方法多针对纯文本信息，缺乏对图像、流程图等多模态信息的处理能力，难以直接应用于电力系统等多模态文档丰富的领域。

\subsubsection{多智能体工作流框架}

多智能体系统的研究最早可追溯至分布式人工智能领域，近年随LLM的突破而获得新的发展动力。AutoGen\cite{wu2023autogen}设计了可对话的智能体编程接口，支持多智能体协作完成代码生成、数据分析等任务。LangGraph基于有向图的计算模型，将智能体工作流抽象为节点（Node）和边（Edge），支持条件分支和循环结构的精确控制。MetaGPT\cite{hong2023metagpt}将软件工程中的角色分工引入多智能体系统，通过标准化输入输出协议规范智能体间通信。

TaskWeaver\cite{taskweaver2024}提出以代码片段为执行单元的智能体编排机制，适合数据分析类任务；CAMEL\cite{li2023camel}探索了基于角色扮演的多智能体对话范式；CrewAI\footnote{CrewAI Inc. CrewAI: Framework for Orchestrating Role-Playing, Autonomous AI Agents. \url{https://github.com/crewAIInc/crewAI}}以"船员"隐喻组织智能体团队，支持任务的层级分解和并行执行。然而，上述框架均需人工预定义工作流拓扑，缺乏根据任务描述和领域知识自动生成最优工作流的能力。

\subsubsection{专业领域智能体应用}

将LLM和多智能体系统应用于电力等专业领域是近期研究的重要趋势。Cheng等\cite{cheng2025gaia}提出面向电力调度的Grid Artificial Intelligent Assistant（GAIA），验证了大语言模型在运行调整、运行监视和黑启动等任务中的应用潜力。在知识图谱方面，变电站设备故障诊断\cite{xiao2022substationkg}和智能电网故障诊断\cite{yao2024smartgridkg}等研究表明，结构化知识表示能够提升故障信息组织和推理分析能力。

然而，现有专业领域智能体方法大多孤立解决知识获取或任务执行的单一问题，尚未形成从多模态知识图谱构建到工作流自动生成的完整闭环方法。本研究正是针对这一空白提出系统性解决方案。

\subsection{论文的主要内容与结构}

\subsubsection{论文的主要内容}

本研究围绕"多模态GraphRAG驱动的智能体工作流自主生成"这一核心命题，开展以下四方面工作：

\textbf{（1）多模态知识图谱构建（算法A）。}设计面向电力领域的多模态文档解析流水线，支持文本、图像和流程图三种模态的实体与关系抽取；构建统一的图节点表示，利用BGE-M3嵌入模型生成高维语义向量，并写入Neo4j图数据库建立向量索引。

\textbf{（2）GraphRAG检索方法（算法B）。}提出以语义相似度检索为种子、BFS子图扩展为骨干的两阶段检索策略；引入LLM重排序机制对候选节点按相关性过滤，并设计结构化上下文格式化方法，将子图信息有效编码为LLM可理解的文本上下文。本研究在 30 条电力基准任务上系统评估了不同检索阶段对下游工作流生成的实际影响，量化了向量种子检索与 BFS 子图扩展的独立贡献边界。

\textbf{（3）子任务-智能体匹配算法（算法C）。}基于智能体特征图（Agent Feature Graph），定义能力嵌入相似度、工具覆盖率和输入输出兼容性三维评分函数；提出工具覆盖率软约束机制和消融实验权重配置，实现对最优智能体的系统化选择。

\textbf{（4）工作流编排与验证（算法D）。}基于LLM推断隐式数据流依赖，构建工作流有向无环图；设计三层迭代验证-修复机制（DAG无环性、IO兼容性、任务完整性，最大迭代次数 5），并实现Mermaid图、BPMN 2.0和LangGraph三种格式的工作流导出。

\subsubsection{论文结构}

本文的组织结构如下：

第一章为绪论，介绍研究背景、问题挑战、国内外现状及论文结构。

第二章为相关工作，系统综述多模态RAG、知识图谱构建和多智能体框架等领域的核心技术进展，分析与本研究的关联与区别。

第三章为方法论，详细阐述多模态知识图谱构建（算法A）和GraphRAG检索（算法B）的技术设计，包括节点表示、向量索引和子图扩展等关键步骤。

第四章为智能体匹配与工作流生成，重点介绍子任务-智能体匹配算法（算法C）和工作流DAG编排与验证机制（算法D），并说明系统整体架构。

第五章为实验与分析，基于电力系统三类场景（OPS/KQA/DA）评估所提方法的有效性，与四种基线方法进行对比，并开展消融实验。

第六章为结论与展望，总结本研究的主要贡献，分析现有方法的局限性，并展望未来研究方向。

\subsection{本章小结}

本章介绍了多模态GraphRAG驱动的智能体工作流自动生成研究的背景与意义，分析了当前领域在知识融合、智能体匹配和工作流质量保证三个层面存在的核心挑战，综述了检索增强生成、图结构知识检索、多智能体框架和专业领域应用四个方向的国内外研究现状，并阐明了本文的主要研究内容与章节结构。
